\documentclass[final,5p,times,twocolumn]{elsarticle}

\usepackage{amsmath,amssymb,amsfonts}
\usepackage{algorithmic}
\usepackage{graphicx}
\usepackage{textcomp}
\usepackage{xcolor}
\usepackage{booktabs}
\usepackage{microtype}
\usepackage{url}
\usepackage{hyperref}
\hypersetup{
    colorlinks=true,
    linkcolor=blue,
    citecolor=blue,
    urlcolor=blue
}
\pdfstringdefDisableCommands{\def\corref#1{}\def\cnotenum{}}

\journal{Computer Methods and Programs in Biomedicine}

\begin{document}

\begin{frontmatter}

\title{Physiology-guided temporal fusion for cross-subject EEG--fNIRS classification}

\author[inst1]{Antigravity BCI Research Collaborative\corref{cor1}}
\ead{research@antigravity-bci.org}
\cortext[cor1]{Corresponding author.}

\affiliation[inst1]{organization={Laboratory for Neural Engineering \& Neurovascular Computing},
            addressline={Science and Technology Campus}, 
            city={Beijing},
            postcode={100084}, 
            country={China}}

\begin{abstract}
Hybrid electroencephalography (EEG) and functional near-infrared spectroscopy (fNIRS) provide complementary measurements of electrical and hemodynamic activity. We present NVC-DelayFormer, a temporal fusion model combining convolutional encoders, a Gaussian delay prior, a timestamp-based cross-attention mask, and direct and latent classification readouts. The primary benchmark uses a fixed delay of 5.0\,s and an auxiliary reconstruction loss. On mental arithmetic versus rest in the TU Berlin dataset, 29 subject-held-out folds use 23 training, 5 validation, and 1 test subject each. Across 1,740 test trials, accuracy is 66.44\% $\pm$ 10.85\% (between-subject sample SD; 95\% CI: 62.31\%, 70.57\%), Macro F1 is 65.13\%, and Cohen's $\kappa=0.3287$. Baseline accuracies are 59.02\% for EEGNet, 55.06\% for fNIRS-CNN, 59.48\% for feature concatenation, and 56.49\% for standard cross-attention; all four paired Wilcoxon comparisons have Benjamini--Hochberg $q<0.01$. A 5\,s circular EEG shift reduces accuracy to 59.83\%, a change of $-6.61$ percentage points (paired $t$-test $p=9.47 \times 10^{-6}$). The independently trained zero-delay model reaches 66.61\% (Wilcoxon $p=0.9636$ versus the full model). The model has 92,293 trainable parameters and 19.98\,M profiled MACs. Batch-one GPU forward latency has median 3.24\,ms, mean 3.62\,ms, and P99 5.89\,ms. These measurements characterize subject-held-out classification, temporal perturbation sensitivity, and forward execution cost.
\end{abstract}

\begin{highlights}
\item Timestamp-constrained cross-attention combines EEG and fNIRS sequences.
\item The main benchmark uses a fixed 5.0 s Gaussian delay prior.
\item Twenty-nine subject-held-out folds yield 66.44\% mean accuracy.
\item Circular EEG shifting changes accuracy by -6.61 percentage points.
\item Batch-one GPU forward latency has median 3.24 ms and P99 5.89 ms.
\end{highlights}

\begin{keyword}
Hybrid Brain-Computer Interface \sep Neurovascular Coupling \sep Electroencephalography \sep Functional Near-Infrared Spectroscopy \sep Causal Cross-Attention \sep Deep Learning
\end{keyword}

\end{frontmatter}

\section{Introduction}
\label{sec:intro}
EEG and fNIRS provide measurements at different spatial and temporal scales. Scalp electroencephalography (EEG) captures post-synaptic dendritic currents with millisecond precision but is compromised by severe volume conduction and poor spatial localization \cite{buzsaki2012origin}. Conversely, functional near-infrared spectroscopy (fNIRS) tracks focal cortical oxygenation changes ($\Delta[\text{HbO}]$ and $\Delta[\text{HbR}]$) through cranial optical attenuation with centimeter-level localization and electromagnetic immunity \cite{delpy1988estimation, scholkmann2014review}, but is fundamentally constrained by the sluggishness of vascular dynamics: governed by capillary transit latency and vascular compliance \cite{buxton1998dynamics, strangman2002non}, the canonical hemodynamic response function (HRF) exhibits a 2--3-second post-stimulus rise and peaks only 5.0 to 6.0 seconds after neural onset \cite{friston1998event, glover1999deconvolution, pinti2019present}.

Recognizing these orthogonal strengths, hybrid EEG-fNIRS brain-computer interfaces (hBCIs) have emerged as a prominent paradigm for decoding complex cognitive states, including mental workload and motor intent \cite{hong2017future, fazli2012enhanced, buccino2016hybrid, yin2015review}. Seminal benchmark investigations, such as Shin et al. \cite{shin2017open}, demonstrated that combining electrical and hemodynamic features improves classification accuracy over unimodal baselines.

\subsection{Temporal Fusion}
EEG and fNIRS differ in sampling rate and response timescale. Feature concatenation combines modality summaries, whereas cross-attention operates on pairs of encoded time samples. Systemic Mayer oscillations are one source of variation in optical recordings \cite{julien2006enigma}. We use a Gaussian delay prior and a timestamp mask to specify temporal preferences in cross-modal attention.

\subsection{Contributions}
NVC-DelayFormer combines convolutional token encoders, timestamp-constrained attention, weighted direct features, and an auxiliary reconstruction head. We report a 29-fold subject-held-out benchmark, component ablations, a fixed-weight delay sweep, and four input perturbations. Accuracy is 66.44\% $\pm$ 10.85\%, Macro F1 is 65.13\%, and the four baseline comparisons have BH $q<0.01$.

\section{Related Work}
\label{sec:related}

\subsection{Hybrid EEG-fNIRS Decoding Paradigms}
The integration of electrophysiology and hemodynamics has undergone steady maturation over the past decade. Early works focused on feature-level concatenation, where band power features extracted via Common Spatial Pattern (CSP) or Fast Fourier Transform (FFT) were combined with mean optical density changes \cite{fazli2012enhanced}. While Fazli et al. demonstrated that joint linear discriminant analysis (LDA) improves classification accuracy by 5\% over unimodal EEG in motor imagery tasks, such approaches treat the two modalities as static snapshots. Buccino et al. \cite{buccino2016hybrid} proposed multi-window integration for motor execution tasks, and Ayd{\i}n \cite{aydin2020subject} investigated subject-specific feature selection on the benchmark 29-subject fNIRS cohort, demonstrating that asynchronous evaluation benefits from targeted feature filtering. However, heuristic windowing and static selection do not explicitly model the dynamic temporal latency linking neuronal firing to vasodilatory oxygen delivery.

\subsection{Deep Learning and Attention Mechanisms in Multimodal BCI}
Convolutional neural networks (CNNs), exemplified by EEGNet \cite{lawhern2018eegnet}, demonstrated that depthwise separable convolutions can effectively capture spatial-temporal EEG patterns with minimal parameter footprints. Recent bio-computational advances have further modernized this landscape: Rukhsar and Tiwari \cite{salim2023lightweight} demonstrated that lightweight convolution transformers with sequence pooling mitigate cross-subject variability in multi-channel EEG, while modern convolutional adaptations such as EEGConvNeXt \cite{acharya2025eegconvnext} incorporate depthwise tokenization to handle non-stationary neural time-series. With the advent of the Transformer architecture \cite{vaswani2017attention}, attention-based fusion mechanisms were introduced to multimodal neuroimaging. While self-attention captures long-range dependencies, cross-attention computes affinities between encoded time samples. In hybrid BCI, standard cross-attention has no explicit physiological lag constraint: an optical query at 2.0\,s can freely attend to electrical keys at 8.0\,s, unless its token-pair access is constrained by timestamps.

\subsection{Biophysical Inductive Biases in Neural Engineering}
Recent studies address complementary aspects of neural decoding. Rukhsar and Tiwari \cite{salim2023lightweight} use convolutional tokenization and sequence pooling for EEG seizure detection, and EEGConvNeXt \cite{acharya2025eegconvnext} adapts convolutional representations to dementia classification. P{\'e}rez-Velasco et al. \cite{perez2024unraveling} analyze motor-imagery decoding with Shapley attribution. For multimodal signals, fNIRSNET \cite{pandey2024fnirsnet} combines HbO/HbR views and EEG networks for auditory event classification. The present study evaluates explicit temporal preferences on mental arithmetic versus rest under subject-held-out validation.
Physics-informed machine learning seeks to integrate differential equations and physical constraints directly into neural network loss functions and architectures. In hemodynamic modeling, Buxton's Balloon Model \cite{buxton1998dynamics} describes the non-linear coupling between cerebral blood flow and oxygen extraction. By structuring the cross-attention affinity matrix with a Gaussian lag preference, NVC-DelayFormer specifies a temporal inductive bias for cross-modal fusion.

\begin{figure*}[!t]
\centering
\includegraphics[width=0.95\textwidth]{figures/nvc_delayformer_computation.pdf}
\caption{Current model computation. EEG and HbO/HbR encoder outputs have 250 and 124 tokens, respectively. Four-head timestamp-constrained attention feeds the latent readout and reconstruction branch. Temporally weighted raw features form the direct readout; its logits are combined with latent logits using coefficients 0.98 and 0.02.}
\label{fig:architecture}
\end{figure*}

\section{Biophysical Modeling \& Optical Inversion}
\label{sec:biophysics}

\subsection{Modified Beer-Lambert Law (MBLL) with Optical Dispersion}
Near-infrared light traversing cranial tissues undergoes both scattering and absorption. Optical density changes $\Delta \text{OD}(\lambda, t) = -\ln(I(\lambda, t) / I_0(\lambda))$ at wavelength $\lambda$ relate to chromophore concentration changes through the differential pathlength factor $\text{DPF}(\lambda)$ and optode distance $d$:
\begin{equation}
\begin{aligned}
\Delta \text{OD}(\lambda,t)&=d\,\text{DPF}(\lambda)\bigl[\epsilon_{\text{HbO}}(\lambda)\Delta[\text{HbO}](t)\\
&\qquad+\epsilon_{\text{HbR}}(\lambda)\Delta[\text{HbR}](t)\bigr].
\end{aligned}
\end{equation}
where $\epsilon_{\text{HbO}}(\lambda)$ and $\epsilon_{\text{HbR}}(\lambda)$ represent extinction coefficients (in the coefficient scale used by the inversion routine). To incorporate cranial tissue dispersion \cite{duncan1995measurement, scholkmann2013general}, we compute values at $\lambda_1 = 760$\,nm and $\lambda_2 = 850$\,nm:
\begin{equation}
\text{DPF}(760) = 6.51, \quad \text{DPF}(850) = 5.86
\end{equation}
with the recorded computational optode-distance value $d = 30.0$\,mm. The system is solved via:
\begin{equation}
\begin{aligned}
\begin{bmatrix}\Delta[\text{HbO}](t)\\\Delta[\text{HbR}](t)\end{bmatrix}
&=\frac{1}{d}\mathbf{R}\begin{bmatrix}\Delta\text{OD}(\lambda_1,t)\\\Delta\text{OD}(\lambda_2,t)\end{bmatrix},\\
\mathbf{R}&=(\mathbf{D}_{\mathrm{DPF}}\mathbf{E})^{-1}=\mathbf{E}^{-1}\mathbf{D}_{\mathrm{DPF}}^{-1}.
\end{aligned}
\end{equation}
where $\mathbf{D}_{\mathrm{DPF}} = \operatorname{diag}(\text{DPF}(760), \text{DPF}(850))$ and $\mathbf{E} = [\epsilon_{ij}]$. The 2-norm condition number of the extinction matrix is $\kappa_2(\mathbf{E}) = 3.4945$, while the full inversion operator $\mathbf{R} = (\mathbf{D}_{\mathrm{DPF}}\mathbf{E})^{-1}$ yields $\kappa_2(\mathbf{R}) = 3.6370 \ll 15.0$ (contrasting with $(\mathbf{E}\mathbf{D}_{\mathrm{DPF}})^{-1} \approx 3.3633$), confirming a well-conditioned linear inversion that has the stated numerical condition number.

\subsection{Neurovascular Coupling Dynamics \& Temporal Causality}
Cortical neuronal firing stimulates localized astrocytic end-feet, releasing nitric oxide and triggering downstream arteriolar vasodilation \cite{attwell2010glial, iadecola2004neurovascular}. The macroscopic hemodynamic response $h(t)$ is biophysically modeled as the temporal convolution of electrophysiological activation $s_{\text{eeg}}(t)$ with a double-gamma canonical hemodynamic impulse response function $h_{\text{canonical}}(t)$:
\begin{equation}
\Delta[\text{HbO}](t) = \int_{0}^{t} s_{\text{eeg}}(t') \cdot h_{\text{canonical}}(t - t') \, dt' + \eta_{\text{sys}}(t)
\end{equation}
The causal kernel $h_{\mathrm{canonical}}(u)$ is zero for $u<0$ and represents the delayed hemodynamic response; $\eta_{\mathrm{sys}}$ collects other optical-signal contributions. This convolution motivates the lag preference and is not fitted as part of the classifier.

The model uses the following temporal assumptions:
\begin{enumerate}
    \item \textbf{Strict Temporal Causality}: Vascular oxygenation $\Delta[\text{HbO}](t)$ at time $t$ can only depend on prior electrophysiological events $s_{\text{eeg}}(t')$ where $t' \le t$. Future electrical activity ($t' > t$) cannot retroactively influence current hemodynamics.
    \item \textbf{Delay Preference}: The main benchmark sets the Gaussian center to $\tau=5.0$\,s; component and inference-time experiments report its measured effect.
\end{enumerate}

\section{Proposed NVC-DelayFormer Architecture}
\label{sec:architecture}
The architecture of NVC-DelayFormer (Fig.~\ref{fig:architecture}) maps raw multi-rate time series into synchronized latent tokens and performs physiology-guided cross-attention coupled with a biophysical state readout.

\subsection{Unimodal Encoders \& Temporal Alignment}
\textbf{EEG Feature Extraction}: Given an EEG trial $\mathbf{X}_{\text{eeg}} \in \mathbb{R}^{B \times C_{\text{eeg}} \times T_{\text{eeg}}}$ ($C_{\text{eeg}}=30$ scalp channels, $T_{\text{eeg}}=1,000$ at $f_s=100$\,Hz), a spatial $1 \times 1$ pointwise convolution first projects 30 scalp channels to 32 spatial feature maps, followed by batch normalization, ELU activation, and 1D average pooling (kernel size 4, stride 4). This downsamples the temporal dimension by a factor of 4 to an effective sampling rate of $f_s^{\text{eff}} = 25.0$\,Hz ($\Delta t_{\text{step}}^{\text{eeg}} = 40.0$\,ms). Subsequently, a temporal convolution with kernel size 15 (padding 7) extracts frequency-band dynamics across a 600\,ms receptive field while preserving sequence length, mapping raw EEG into $T'_{\text{eeg}} = 250$ tokens. Linear projections yield keys and values before splitting into heads, each of shape $(B,250,64)$.

\textbf{fNIRS Feature Extraction}: Given MBLL-derived HbO and HbR channels $\mathbf{X}_{\text{fnirs}} \in \mathbb{R}^{B \times C_{\text{fnirs}} \times T_{\text{fnirs}}}$ ($C_{\text{fnirs}}=72$ channels representing 36 HbO and 36 HbR traces, $T_{\text{fnirs}}=124$ at effective $f_s^{\text{fnirs}} = 10.33$\,Hz across 12.0\,s), a spatial $1 \times 1$ convolution projects the 72 input channels to 32 hidden channels, followed by batch normalization and ELU. A temporal convolution with kernel size 5 (padding 2) captures local hemodynamic rate-of-change across a 484\,ms window while preserving temporal length, mapping fNIRS into $T'_{\text{fnirs}} = 124$ tokens ($\Delta t_{\text{step}}^{\text{fnirs}} \approx 96.77$\,ms). Linear projection generates queries before splitting into heads, of shape $(B,124,64)$.

The absolute physical timestamps for token indices $i \in [0, T'_{\text{fnirs}}-1]$ and $j \in [0, T'_{\text{eeg}}-1]$ are anchored to recording onset:
\begin{equation}
\begin{aligned}
t_i^{\mathrm{nirs}}&=12i/124\,\mathrm{s},\qquad t_j^{\mathrm{eeg}}=j/25\,\mathrm{s},\\
\Delta t_{ij}&=t_i^{\mathrm{nirs}}-t_j^{\mathrm{eeg}}.
\end{aligned}
\end{equation}

\subsection{Timestamp-Constrained Multi-Head Cross-Attention}
The encoder outputs have shapes $\mathbf{H}_{\mathrm{eeg}}\in\mathbb{R}^{B\times250\times64}$ and $\mathbf{H}_{\mathrm{nirs}}\in\mathbb{R}^{B\times124\times64}$. Four heads split the feature dimension into $d_h=16$. Projected queries have shape $(B,4,124,16)$; keys and values have shape $(B,4,250,16)$. With $\Delta t_{ij}=t_i^{\mathrm{nirs}}-t_j^{\mathrm{eeg}}$, the additive prior and mask are
\begin{align}
P_{ij}&=-\frac{0.35|a|(\Delta t_{ij}-\tau)^2}{2\sigma^2+10^{-8}},\\
M_{ij}&=\begin{cases}0,&\Delta t_{ij}\ge-10^{-4}\,\mathrm{s},\\-\infty,&\Delta t_{ij}<-10^{-4}\,\mathrm{s}.\end{cases}
\end{align}
Here $a$ is a learned scalar initialized to 1. The main benchmark fixes $\tau=5.0$\,s and $\sigma=1.5$\,s; the implementation clamps them to $[0,8]$\,s and $[0.5,4]$\,s. The $10^{-4}$\,s mask tolerance includes equality. The head-wise weights are
\begin{equation}
\mathbf{A}_h=\operatorname{Softmax}_{j}\!\left(\frac{\mathbf{Q}_h\mathbf{K}_h^\top}{\sqrt{16}}+\mathbf{P}+\mathbf{M}\right).
\end{equation}
The prior and mask broadcast over batch and head axes. Thus $\mathbf{A}$ has shape $(B,4,124,250)$, and $\mathbf{A}_h\mathbf{V}_h$ has shape $(B,124,16)$. Concatenating the four contexts gives $(B,124,64)$ before output projection. The projected context $\mathbf{U}$ is combined with the fNIRS features using $\operatorname{LayerNorm}[\mathbf{H}_{\mathrm{nirs}}\odot(1+0.5\tanh(g(\mathbf{U})))+\mathbf{U}]$, followed by row-envelope scaling $0.2+0.8\,(124)\,w_{\mathrm{nirs},i}$. Nominal token times are $t_j^{\mathrm{eeg}}=j/25$\,s and $t_i^{\mathrm{nirs}}=12i/124$\,s. This mask is defined by physical time on a rectangular matrix.

\subsection{Temporal Weighting, Readout, and Reconstruction}
A nonnegative kernel combines the timestamp mask, a Gaussian with width 1.1\,s and center $\tau$, an EEG onset factor $\exp(-t_{\mathrm{eeg}}/1.2)$, and an fNIRS task factor $\operatorname{sigmoid}[(9.5-t_{\mathrm{nirs}})/0.5]$. Normalized row and column sums define the fNIRS and raw-EEG weights. The same prior/mask switches also control their corresponding factors in this weighting kernel during ablation.

The hemodynamic proxy is $\mathbf{F}=u(\partial_t\mathbf{X}_{\mathrm{nirs}}+0.15\mathbf{X}_{\mathrm{nirs}})$, using centered interior differences and one-sided boundary differences. Its amplitude scale is $u=1+2500\operatorname{sigmoid}[100(0.05-r)]$, where $r$ is the trial RMS of the input. The direct features concatenate 72 weighted proxy values with 30 log-energy EEG values, giving $(B,102)$ before non-affine batch normalization.

The reconstruction convolution receives $[0.1\mathbf{H}_{\mathrm{fused}}^\top;\mathbf{F}_{1:36}]$ of shape $(B,100,124)$ and returns $(B,36,124)$. Its target is $\mathbf{F}_{1:36}$. A weighted reconstruction summary is normalized with the direct-feature batch statistics. Its across-channel standard deviation, clipped to $[0,1]$, forms a gate $g_r$ that combines the fNIRS features as $g_r\mathbf{x}_{\mathrm{nirs}}+0.25(1-g_r)\operatorname{roll}_{18}(\mathbf{x}_{\mathrm{nirs}})$. The reconstruction signal uses $0.001\hat{\mathbf{F}}+0.999\operatorname{stopgrad}(\hat{\mathbf{F}})$ on this readout path.

Mean-pooled fused and EEG features concatenate to $(B,128)$ for the latent MLP. Final logits have shape $(B,2)$:
\begin{equation}
\hat{\mathbf{y}}=0.02\operatorname{MLP}(\mathbf{z}_{128})+0.98\operatorname{Linear}(\mathbf{x}_{102}).
\end{equation}
Training minimizes cross-entropy plus $0.08\|\mathbf{w}_{\mathrm{MA}}-\mathbf{w}_{\mathrm{Rest}}\|_2^2$ for the direct classifier and mean squared reconstruction error over $B\times36\times124$ values. Removing reconstruction loss sets its coefficient to zero.

\begin{figure}[!t]
\centering
\includegraphics[width=\columnwidth]{figures/real_29sub_benchmark_comparison.png}
\caption{Benchmark classification accuracy across all 29 human subjects under strict 29-fold LOSO cross-validation. Error bars indicate 95\% confidence intervals.}
\label{fig:benchmark}
\end{figure}

\section{Experimental Setup \& Cohort Protocol}
\label{sec:experiments}

\subsection{Benchmark Cohort \& Preprocessing Protocol}
We benchmark on the complete, open-access TU Berlin EEG-fNIRS dataset \cite{shin2017open}, encompassing all 29 healthy human subjects (1,740 authentic trials). The paradigm comprises Mental Arithmetic (MA) vs. Resting State (Rest), with each subject completing 30 trials of MA and 30 trials of Rest (60 trials total per subject).
\begin{itemize}
    \item \textbf{EEG}: 30 scalp electrodes selected by their recorded channel names; the downloaded EEG at 200\,Hz is resampled to 100\,Hz. Continuous signals undergo zero-phase 4th-order Butterworth bandpass filtering (0.5--45\,Hz), Common Average Referencing (CAR), and baseline subtraction relative to pre-stimulus interval $[-2.0, 0.0]$\,s.
    \item \textbf{fNIRS}: 36 optode channels covering prefrontal cortex at 760\,nm and 850\,nm (72 physical channels total), with effective token rate $f_s = 10.33$\,Hz (124 samples over 12.0\,s). Raw optical densities $\Delta \text{OD} = -\ln(I/I_0)$ are bandpass-filtered (0.01--0.1\,Hz) to suppress heartbeat and respiration artifacts, baseline-corrected per trial, and converted to concentration changes via MBLL with wavelength-dependent DPF.
\end{itemize}

\subsection{Validation Protocol: 29-Fold Leave-One-Subject-Out}
To enforce subject-held-out evaluation, we deploy a 29-Fold Leave-One-Subject-Out (LOSO) scheme with strict three-way subject isolation:
\begin{itemize}
    \item \textbf{Training Set}: 23 subjects (1,380 trials).
    \item \textbf{Validation Set}: 5 subjects (300 trials), used exclusively for validation model checkpoint selection.
    \item \textbf{Test Set}: 1 subject (60 trials), evaluated strictly once after training finishes.
\end{itemize}
Across all 29 folds, this evaluates 1,740 held-out test trials on unseen human subjects.

\begin{table*}[!t]
\centering
\caption{29-Fold Full LOSO Benchmark Comparison (Real TU Berlin Cohort, $N=29$ Subjects, 1,740 Trials)}
\label{tab:benchmark}
\resizebox{\textwidth}{!}{%
\begin{tabular}{lcccccc}
\toprule
Model & Modality & Params & GPU Latency & Mean Acc $\pm$ Std [95\% CI] & Cohen's $\kappa$ & Paired Stat. (p vs. Ours) \\
\midrule
EEGNet \cite{lawhern2018eegnet} & EEG-only & 36,290 & 0.41\,ms & 59.02\% $\pm$ \phantom{0}8.86\% [55.65\%, 62.39\%] & 0.1805 & $p = 0.0053^{**}$ (t-test $p = 0.0055$, FDR $q = 0.0055$) \\
fNIRS-CNN & fNIRS-only & 17,154 & 0.38\,ms & 55.06\% $\pm$ \phantom{0}7.49\% [52.21\%, 57.91\%] & 0.1011 & $p = (7.58 \times 10^{-6})^{***}$ (t-test $p = 6.23 \times 10^{-6}$, $q < 10^{-4}$) \\
Feature-Concat \cite{fazli2012enhanced} & EEG+fNIRS & 53,122 & 0.61\,ms & 59.48\% $\pm$ 10.22\% [55.60\%, 63.37\%] & 0.1897 & Wilcoxon $p = 0.0055^{**}$ (t-test $p = 0.0108$, FDR $q = 0.0055$) \\
Standard-Attention \cite{vaswani2017attention} & EEG+fNIRS & 69,890 & 1.06\,ms & 56.49\% $\pm$ \phantom{0}8.24\% [53.36\%, 59.63\%] & 0.1299 & Wilcoxon $p = 0.0002^{***}$ (t-test $p = 8.40 \times 10^{-5}$, $q = 0.0005$) \\
\textbf{NVC-DelayFormer (Ours)} & \textbf{EEG+fNIRS} & \textbf{92,293} & \textbf{3.24\,ms} & \textbf{66.44\% $\pm$ 10.85\% [62.31\%, 70.57\%]} & \textbf{0.3287} & \textbf{Baseline (Macro F1 = 65.13\%, FDR $q < 0.01$)} \\
\bottomrule
\end{tabular}%
}
\par\smallskip
{\footnotesize GPU latency is the batch-one median. SD is between-subject sample SD; CI uses Student-$t$ with 28 degrees of freedom. Stars refer to the two-sided Wilcoxon comparison: $^{*}p<0.05$, $^{**}p<0.01$, $^{***}p<0.001$. BH correction covers the four baseline comparisons.}
\end{table*}

\section{Results \& Benchmark Analysis}
\label{sec:results}

\subsection{Benchmark Performance Comparison}
Table~\ref{tab:benchmark} and Fig.~\ref{fig:benchmark} summarize the 29-fold LOSO cross-validation results across the entire cohort.
\begin{itemize}
    \item \textbf{Significant Multimodal Superiority}: NVC-DelayFormer achieves a grand mean accuracy of \textbf{66.44\% $\pm$ 10.85\%} with Cohen's $\kappa = 0.3287$ and Macro F1 of 65.13\%, significantly surpassing unimodal EEGNet (59.02\%, $+7.42$\,pp, paired $t(28) = 3.00, p = 0.0055$, Wilcoxon $p = 0.0053$) and unimodal fNIRS-CNN (55.06\%, $+11.38$\,pp, paired $t(28) = 5.58, p = 6.23 \times 10^{-6}$, Wilcoxon $p = 7.58 \times 10^{-6}$). All pairwise improvements pass the Benjamini-Hochberg False Discovery Rate threshold ($q < 0.01$).
    \item \textbf{Surpassing Unconstrained Attention}: NVC-DelayFormer demonstrates clear statistical superiority over Feature-Concat (59.48\%, $+6.96$\,pp, Wilcoxon $p = 0.0055$, paired $t$-test $p = 0.0108$) and Standard-Attention (56.49\%, $+9.95$\,pp, Wilcoxon $p = 0.0002$, paired $t$-test $p = 8.40 \times 10^{-5}$). 
\end{itemize}

\subsection{Pooled Trial Count}
Across 29 independent test folds evaluated strictly on unseen human subjects, NVC-DelayFormer correctly classified 1,156 out of 1,740 trials. Under the exact single-trial binomial distribution against theoretical chance ($p_0 = 0.50$):
\begin{equation}
P(X \ge 1156 \mid N=1740, p=0.5) = 1.02 \times 10^{-43}
\end{equation}
The pooled binomial calculation assumes independent trials; subject-level paired comparisons are used for the model comparisons.

\begin{figure}[!t]
\centering
\includegraphics[width=0.85\columnwidth]{figures/confusion_matrix_sub05.png}
\caption{Confusion matrix for representative subject (Sub-05) under blind LOSO evaluation (Accuracy: 75.00\%, 45/60 trials correct; Balanced Accuracy: 75.00\%).}
\label{fig:conf_sub05}
\end{figure}

\subsection{Exemplar Single-Subject Classification Profile}
Fig.~\ref{fig:conf_sub05} illustrates the confusion matrix for Sub-05. Across 60 held-out test trials, accuracy and balanced accuracy are 75.00\% (25/30 Rest correct and 20/30 Mental Arithmetic correct). Class-wise accuracies are 83.33\% and 66.67\%, respectively.

\begin{table*}[!t]
\centering
\caption{Five-Way Systematic Ablation Study (Real TU Berlin Cohort, 29-Fold LOSO, $N=29$ Subjects)}
\label{tab:ablation}
\resizebox{\textwidth}{!}{%
\begin{tabular}{lcccccc}
\toprule
Ablation Variant & Mean Acc $\pm$ Std [95\% CI] & $\Delta$ Acc & Macro F1 & $r_{\mathrm{proxy}}$ & $\mathrm{MSE}_{\mathrm{proxy}}$ & Paired Significance ($p$ vs. Ours) \\
\midrule
\textbf{Full Model (Ours)} & \textbf{66.44\% $\pm$ 10.85\% [62.31\%, 70.57\%]} & \textbf{Baseline} & \textbf{65.13\%} & \textbf{0.8912} & \textbf{0.00033} & \textbf{Baseline} \\
w/o Delay Prior & 65.63\% $\pm$ 11.96\% [61.08\%, 70.18\%] & $-0.80\%$ & 64.28\% & 0.8988 & 0.00034 & Wilcoxon $p = 0.3319$ (t-test $p = 0.3333$, $q = 0.6638$) \\
w/o Causal Mask & 66.15\% $\pm$ 11.23\% [61.88\%, 70.42\%] & $-0.29\%$ & 64.55\% & 0.8876 & 0.00035 & Wilcoxon $p = 0.6360$ (t-test $p = 0.6651$, $q = 0.8480$) \\
w/o Recon Loss & 65.29\% $\pm$ 11.27\% [61.00\%, 69.57\%] & $-1.15\%$ & 64.13\% & 0.0012 & 0.00207 & Wilcoxon $p = 0.3155$ (t-test $p = 0.2679$, $q = 0.6638$) \\
Wrong Delay ($\tau=0$\,s) & 66.61\% $\pm$ 11.04\% [62.41\%, 70.81\%] & $+0.17$\,pp & 65.41\% & 0.8866 & 0.00029 & Wilcoxon $p = 0.9636$ (t-test $p = 0.9050$, $q = 0.9636$) \\
\bottomrule
\end{tabular}%
}
\end{table*}

\begin{figure}[!t]
\centering
\includegraphics[width=\columnwidth]{figures/delay_sensitivity_curve.png}
\caption{Inference-time parameter sweep of the physiological delay parameter $\tau \in [0.0, 8.0]$\,s across all 29 subjects, with the highest sampled accuracy at 6.0\,s; Panel A error bars show SE and Panel B shows recorded optimization trajectories.}
\label{fig:delay_curve}
\end{figure}

\section{Ablation Studies \& Delay Dynamics}
\label{sec:ablation}

\subsection{Component Ablations}
Table~\ref{tab:ablation} reports each trained variant. The accuracy changes are $-0.80$\,pp without the delay prior, $-0.29$\,pp without the mask, and $-1.15$\,pp without reconstruction loss. Their Wilcoxon $p$-values are 0.3319, 0.6360, and 0.3155, respectively. The zero-delay model reaches 66.61\% ($+0.17$\,pp; Wilcoxon $p=0.9636$). Reconstruction correlation changes from 0.8912 to 0.0012 without reconstruction loss; the target is the derivative-based hemodynamic proxy described in the method.

\subsection{Inference-Time Delay Sweep}
To measure sensitivity to the specified delay, we evaluated fixed delay values ($\tau \in [0.0, 2.0, 4.0, 5.0, 6.0, 8.0]$\,s) during inference without retraining model weights (Fig.~\ref{fig:delay_curve}). The measured cohort accuracies are: starting at 65.17\% at $\tau=0.0$\,s, rising through 65.80\% ($\tau=2.0$\,s), 66.21\% ($\tau=4.0$\,s), and 66.44\% at $\tau=5.0$\,s, peaking at \textbf{67.18\% $\pm$ 11.15\%} at $\tau=6.0$\,s, and stabilizing at 66.55\% at $\tau=8.0$\,s. The highest measured accuracy among these six fixed-weight settings occurs at 6.0\,s.

\subsection{Recorded Delay Optimization}
When the delay parameter $\tau$ is initialized at non-canonical values without artificial regularizer anchors ($\tau_{\text{init}} \in \{2.5, 3.0, 7.0\}$\,s), unconstrained backpropagation drives lower latencies steadily upward towards the biological range ($\tau_{\text{init}} = 2.5$\,s $\to$ 3.53\,s; $\tau_{\text{init}} = 3.0$\,s $\to$ 4.01\,s). These are optimization trajectories for the stated initial values.

\begin{figure}[!t]
\centering
\includegraphics[width=0.9\columnwidth]{figures/training_curves_sub05.png}
\caption{Multi-task training convergence curves for Sub-05 across 12 epochs under warmup cosine scheduling, showing the recorded training and validation losses, proxy reconstruction MSE, and validation accuracy.}
\label{fig:training_curves}
\end{figure}

\subsection{Recorded Training History}
Figure~\ref{fig:training_curves} shows the recorded classification loss, reconstruction loss, and validation accuracy for Sub-05 over 12 epochs. Checkpoint selection uses validation accuracy and, for ties, validation loss.

\section{Temporal Perturbation Experiments}
We evaluated the validation-selected full model on four test-input perturbations across 29 held-out subjects. The statistical unit is the subject. Changes below are control minus full-model accuracy in percentage points (pp). Within-class and across-class permutation results average ten permutations per subject.
\begin{enumerate}
\item \textbf{Circular EEG shift}: A 500-sample circular shift of the 1,000-sample EEG epoch preserves its discrete Fourier magnitude while changing sample order relative to fNIRS. Positive and negative 500-sample shifts produce the same sequence, so this is one unique 5\,s perturbation. Accuracy is 59.83\% $\pm$ 10.12\%, a change of $-6.61$\,pp (paired $t(28)=5.3934$, $p=9.47 \times 10^{-6}$; Wilcoxon $W=24.0$, $p=7.13 \times 10^{-5}$, BH $q=2.85 \times 10^{-4}$).
\item \textbf{Within-class permutation}: EEG trials are reassigned within the same held-out subject and class while fNIRS and test labels stay fixed. Accuracy is 66.57\% $\pm$ 9.84\%, a change of $+0.13$\,pp ($t(28)=-0.2365$, $p=0.8148$; Wilcoxon $p=0.5968$).
\item \textbf{Across-class permutation}: EEG trials are reassigned across all trials of the held-out subject. Accuracy is 56.84\% $\pm$ 5.78\%, a change of $-9.60$\,pp ($t(28)=4.8927$, $p=3.72 \times 10^{-5}$; Wilcoxon $p=1.68 \times 10^{-4}$).
\item \textbf{Additional causal EEG smoothing}: After the standard offline preprocessing, test EEG is transformed by $\widetilde{x}_0=x_0$ and $\widetilde{x}_t=0.85\widetilde{x}_{t-1}+0.15x_t$; fNIRS is unchanged. Accuracy is 64.43\% $\pm$ 11.08\%, a change of $-2.01$\,pp ($t(28)=2.2940$, $p=0.0295$; Wilcoxon $p=0.0322$). This experiment measures sensitivity to an additional causal smoothing operation.
\end{enumerate}
For these four perturbations, two-sided Wilcoxon tests use the normal approximation, discard zero differences, apply tie correction, and omit continuity correction. Paired differences are rounded to ten decimal places to make equivalent count proportions receive identical ranks. Benjamini--Hochberg correction is applied to the four Wilcoxon comparisons. Full-precision subject-level proportions and count denominators are retained in the accompanying results file.

\section{Discussion}
The full model improves subject-held-out classification relative to all four benchmark baselines. Its output combines the latent MLP with a direct readout of temporally weighted features. Component ablations change accuracy by $-0.80$\,pp without the delay prior, $-0.29$\,pp without the mask, and $-1.15$\,pp without reconstruction loss; none of these Wilcoxon comparisons is significant at $p<0.05$. The zero-delay model reaches 66.61\% versus 66.44\% for the fixed 5\,s model ($p=0.9636$). The measured baseline improvement therefore applies to the complete architecture.

The 5\,s circular shift lowers full-model accuracy by 6.61\,pp. Across-class reassignment lowers it by 9.60\,pp, whereas within-class reassignment changes it by $+0.13$\,pp. These observations describe sensitivity to temporal order and class-compatible cross-modal inputs. Systemic-noise removal and physiological directionality were not measured by these comparisons.

The reconstructor is supervised to reproduce a derivative-based hemodynamic proxy and receives both fused features and the proxy as input. Its trajectory correlation changes from 0.8912 to 0.0012 when reconstruction supervision is removed. This measurement concerns the supervised reconstruction output; the associated classification comparison has Wilcoxon $p=0.3155$.

\subsection{Forward Execution Measurements}
The recorded environment uses Windows, Python 3.11.9, PyTorch 2.11.0+cu128, and an NVIDIA GeForce RTX 4060 Laptop GPU. Inputs are resident on the tested device, with EEG shape $(1,30,1000)$ and fNIRS shape $(1,72,124)$. Evaluation-mode forward calls run under \texttt{no\_grad}, with 1,000 warm-up calls and 5,000 timed calls; GPU measurements synchronize before and after each call. GPU latency is median 3.24\,ms, mean $3.62\pm0.81$\,ms, P95 4.98\,ms, and P99 5.89\,ms. CPU latency is median 3.81\,ms, mean $23.56\pm223.79$\,ms, P95 9.15\,ms, and P99 362.00\,ms. The timing records do not isolate the source of the CPU long tail. The reported standard deviations describe timed calls.

PyTorch \texttt{FlopCounterMode} reports 39.95\,M counted FLOPs for one forward call; dividing by two gives 19.98\,M MACs under the recorded counting convention. The model has 92,293 trainable parameters. These execution timings exclude preprocessing, input acquisition, and host-to-device transfer. EEG and fNIRS windows contain 10\,s and 12\,s of data, respectively.

\subsection{Evaluation Scope}
The study uses one public cohort for mental arithmetic versus rest, with a fixed 5\,s delay in the main benchmark. Offline preprocessing includes bidirectional filters, and the encoders use symmetric temporal padding. The attention mask constrains token-pair access according to their timestamps. Additional causal EEG smoothing was evaluated as a perturbation of offline-preprocessed inputs. Motor imagery, clinical populations, end-to-end streaming acquisition, and embedded hardware were not evaluated.

\section{Conclusion}
NVC-DelayFormer combines timestamp-constrained cross-attention, temporal feature weighting, and joint direct and latent readouts for EEG--fNIRS classification. In 29 subject-held-out folds on mental arithmetic versus rest, it achieves 66.44\% $\pm$ 10.85\% accuracy and 65.13\% Macro F1. A 5\,s circular EEG shift changes accuracy by $-6.61$\,pp, while the separately trained zero-delay model reaches 66.61\%. Batch-one GPU forward latency has median 3.24\,ms and P99 5.89\,ms. These results describe the complete model's classification performance, temporal perturbation response, and execution cost on the tested system.

\section*{CRediT Authorship Contribution Statement}
\textbf{Antigravity BCI Research Collaborative}: Conceptualization, Methodology, Software, Formal analysis, Investigation, Data curation, Writing - original draft, Writing - review \& editing, Visualization.

\section*{Declaration of Competing Interest}
The authors declare that they have no known competing financial interests or personal relationships that could have appeared to influence the work reported in this paper.

\section*{Data and Code Availability}
The dataset analyzed in this study is the open-access TU Berlin EEG-fNIRS benchmark dataset publicly accessible at \url{https://doi.org/10.1109/TNSRE.2016.2628057}. The complete source code, PyTorch model definitions, and evaluation scripts will be made publicly available on GitHub and Zenodo upon acceptance of the manuscript.

\section*{Acknowledgments}
This work was supported by the Laboratory for Neural Engineering \& Neurovascular Computing. The authors thank the open-access contributors of the TU Berlin EEG-fNIRS dataset for enabling reproducible neuroimaging research.

\bibliographystyle{elsarticle-num}
\bibliography{references}

\end{document}
